\QHeading{第1题　用简便方法计算下面各题。}
\par\noindent


\Problem{（1）$101\times23$}
\par


\MethodTitle{方法一：简便计算}
\Answer{{\Formula{\begin{aligned}[t]&=(100+1)\times23\\&=100\times23+1\times23\\&=2300+23\\&=2323\end{aligned}}}}

\MethodBlock{方法二：按数位拆分计算}{把一个因数拆成整十、整百数与余下部分之和。}{\begin{aligned}[t]&\phantom{={}}101\times 23\\
&=101\times (20+3)\\
&=101\times 20+101\times 3\\
&=2020+303\\
&=2323\end{aligned}}
\Problem{（2）$50\times35\times40$}

\MethodTitle{方法一：简便计算}
\Answer{{\Formula{\begin{aligned}[t]&=50\times40\times35\\&=(50\times40)\times35\\&=2000\times35\\&=70000\end{aligned}}}}

\MethodBlock{方法二：结合因数，先凑整}{用乘法交换律、结合律，先算容易凑整的两个因数。}{\begin{aligned}[t]&\phantom{={}}50\times 35\times 40\\
&=(35\times 40)\times 50\\
&=1400\times 50\\
&=70000\end{aligned}}

\MethodBlock{方法三：按运算顺序直接计算（对照）}{不先用运算律变形，按括号、乘除、加减的顺序计算。题目要求简算时，优先使用上面的简便解法。}{\begin{aligned}[t]&\phantom{={}}50\times 35\times 40\\
&=1750\times 40\\
&=70000\end{aligned}}
\Problem{（3）$25\times39\times4$}

\MethodTitle{方法一：简便计算}
\Answer{{\Formula{\begin{aligned}[t]&=25\times4\times39\\&=(25\times4)\times39\\&=100\times39\\&=3900\end{aligned}}}}

\MethodBlock{方法二：按运算顺序直接计算（对照）}{不先用运算律变形，按括号、乘除、加减的顺序计算。题目要求简算时，优先使用上面的简便解法。}{\begin{aligned}[t]&\phantom{={}}25\times 39\times 4\\
&=975\times 4\\
&=3900\end{aligned}}
\Problem{（4）$237\times12-37\times12$}

\MethodTitle{方法一：简便计算}
\Answer{{\Formula{\begin{aligned}[t]&=(237-37)\times12\\&=200\times12\\&=2400\end{aligned}}}}

\MethodBlock{方法二：按运算顺序直接计算（对照）}{不先用运算律变形，按括号、乘除、加减的顺序计算。题目要求简算时，优先使用上面的简便解法。}{\begin{aligned}[t]&\phantom{={}}237\times 12-37\times 12\\
&=2844-444\\
&=2400\end{aligned}}
\Problem{（5）$400\div8\div25$}

\MethodTitle{方法一：简便计算}
\Answer{{\Formula{\begin{aligned}[t]&=400\div(8\times25)\\&=400\div200\\&=2\end{aligned}}}}

\MethodBlock{方法二：先处理括号或乘除，再拆分除数，连续相除}{}{\begin{aligned}[t]&\phantom{={}}400\div 8\div 25\\
&=50\div 25\\
&=50\div (5\times 5)\\
&=50\div 5\div 5\\
&=10\div 5\\
&=2\end{aligned}}

\MethodBlock{方法三：先处理括号或乘除，再利用商不变的性质}{被除数和除数同时乘同一个非零数，商不变。}{\begin{aligned}[t]&\phantom{={}}400\div 8\div 25\\
&=50\div 25\\
&=(50\times 2)\div (25\times 2)\\
&=100\div 50\\
&=2\end{aligned}}
\Problem{（6）$840\div12-600\div12$}

\MethodTitle{方法一：简便计算}
\Answer{{\Formula{\begin{aligned}[t]&=(840-600)\div12\\&=240\div12\\&=20\end{aligned}}}}

\MethodBlock{方法二：按运算顺序直接计算（对照）}{不先用运算律变形，按括号、乘除、加减的顺序计算。题目要求简算时，优先使用上面的简便解法。}{\begin{aligned}[t]&\phantom{={}}840\div 12-600\div 12\\
&=70-50\\
&=20\end{aligned}}

\QHeading{第2题　填上适当的分数或小数。}

\PlainItem{（1）}\par\noindent\begin{adjustbox}{max width=\linewidth,center}\begin{tikzpicture}[x=1cm,y=1cm,font=\normalsize,draw=black,text=black]
\draw[-{Stealth[length=4pt]},line width=.6pt](-.1,0)--(7.25,0);
\foreach \i in {0,...,10}\draw ({\i*.68},0)--({\i*.68},.1);
\node[above,inner sep=1pt] at(0,.12){$0$};
\node[above,inner sep=1pt] at(6.8,.12){$1$};
\node[above,inner sep=1pt] at(0.6800,.13){$0.1$};
\node[draw=answerred,text=answerred,inner sep=3pt] at(0.6800,1.2){$\displaystyle \frac{1}{10}$};
\node[above,inner sep=1pt] at(4.7600,.13){$\frac{7}{10}$};
\node[draw=answerred,text=answerred,inner sep=3pt] at(4.7600,1.2){$\displaystyle 0.7$};
\path[use as bounding box](-.2,-.08) rectangle (7.4,1.65);
\end{tikzpicture}\end{adjustbox}\par
\PlainItem{（2）}\par\noindent\begin{adjustbox}{max width=\linewidth,center}\begin{tikzpicture}[x=1cm,y=1cm,font=\normalsize,draw=black,text=black]
\draw[-{Stealth[length=4pt]},line width=.6pt](-.1,0)--(7.25,0);
\foreach \i in {0,...,10}\draw ({\i*.68},0)--({\i*.68},.1);
\node[above,inner sep=1pt] at(0,.12){$0$};
\node[above,inner sep=1pt] at(6.8,.12){$0.1$};
\node[above,inner sep=1pt] at(2.0400,.13){$\frac{3}{100}$};
\node[draw=answerred,text=answerred,inner sep=3pt] at(2.0400,1.2){$\displaystyle 0.03$};
\node[above,inner sep=1pt] at(3.4000,.13){$0.05$};
\node[draw=answerred,text=answerred,inner sep=3pt] at(3.4000,1.2){$\displaystyle \frac{5}{100}$};
\path[use as bounding box](-.2,-.08) rectangle (7.4,1.65);
\end{tikzpicture}\end{adjustbox}\par
\par\noindent

\QHeading{第3题　解决问题。}

\Problem{① 在一个直角三角形中，一个锐角是$35^\circ$，另一个锐角是多少度？}
\par


\MethodTitle{方法一}
\Answer{{\Formula{\begin{aligned}[t]&\phantom{={}}180^\circ-90^\circ-35^\circ\\&=90^\circ-35^\circ\\&=55^\circ\end{aligned}}}\par 答：另一个锐角是$55^\circ$。}
\Answer{$55^\circ$}
\Answer{$35^\circ$}

\MethodTitle{方法二：直接利用直角三角形两锐角和}
\Answer{直角占90度，两个锐角共90度。}
\Answer{\Formula{\begin{aligned}[t]&\phantom{={}}90-35\\
&=55\text{（度）}\end{aligned}}}
\Answer{答：另一个锐角是55度。}
\Problem{② 校园里有一个直角三角形的花坛，它的较大锐角是较小锐角的2倍。你能算出这个花坛每个角的度数吗？}

\MethodTitle{方法一}
\Answer{较小角：{\Formula{\begin{aligned}[t]&\phantom{={}}(180^\circ-90^\circ)\div(1+2)\\&=90^\circ\div3\\&=30^\circ\end{aligned}}}\par 较大角：$30^\circ\times2=60^\circ$。\par 答：三个角为$30^\circ$、$60^\circ$、$90^\circ$。}
\Answer{$60^\circ$}
\Answer{$30^\circ$}

\MethodTitle{方法二：按份数把全部180度分成6份}
\Answer{小锐角1份，大锐角2份；直角等于这两个锐角的和，是3份；总共6份。}
\Answer{\textbf{每份度数}\par \Formula{\begin{aligned}[t]&\phantom{={}}180\div (1+2+3)\\
&=180\div (3+3)\\
&=180\div 6\\
&=30\text{（度）}\end{aligned}}}
\Answer{\textbf{大锐角}\par \Formula{\begin{aligned}[t]&\phantom{={}}30\times 2\\
&=60\text{（度）}\end{aligned}}}
\Answer{\textbf{直角}\par \Formula{\begin{aligned}[t]&\phantom{={}}30\times 3\\
&=90\text{（度）}\end{aligned}}}
\Answer{答：三个角分别是30度、60度、90度。}
